\studySubsection{practice-calc-05}{练习库：高等数学第 5 讲 一元函数微分学的应用（一）——几何应用}

\problemAnchor{MATH1-CALC-0019}
\begin{problemBox}
\begin{problemMeta}
\textbf{题目编号：} \problemRef{MATH1-CALC-0019} \quad
\textbf{答案：} \answerRef{MATH1-CALC-0019} \quad
\textbf{来源：} Codex 原创 / K070 深度讲义配套练习（非真题）
\end{problemMeta}
\begin{enumerate}
  \item[\textbf{A}] 求 \(f(x)=xe^{-x}\) 的最大单调区间。
  \item[\textbf{B}] 求 \(f(x)=|x-1|+\lvert x+1\rvert\) 的单调区间，并说明常值区间怎样表述。
  \item[\textbf{C}] 讨论 \(a\in\mathbb R\) 时 \(f_a(x)=x+ae^{-x}\) 在 \(\mathbb R\) 上的单调性。
  \item[\textbf{M}] 证明方程 \(xe^x=1\) 在 \((0,\infty)\) 内恰有一个实根。
\end{enumerate}
\end{problemBox}

\problemAnchor{MATH1-CALC-0020}
\begin{problemBox}
\begin{problemMeta}
\textbf{题目编号：} \problemRef{MATH1-CALC-0020} \quad
\textbf{答案：} \answerRef{MATH1-CALC-0020} \quad
\textbf{来源：} Codex 原创 / K071 深度讲义配套练习（非真题）
\end{problemMeta}
\begin{enumerate}
  \item[\textbf{A}] 判断 \(x=0\) 是否为 \(f(x)=x^4(1+x)\) 的极值点，并说明二阶导判据为何不能直接给出结论。
  \item[\textbf{B}] 求 \(f(x)=x^2e^{-x}\) 的全部局部极值。
  \item[\textbf{C}] 求 \(f(x)=|x^2-1|\) 的全部局部极值点与极值。
  \item[\textbf{M}] 讨论 \(a\in\mathbb R\) 时 \(f(x)=x^4-2ax^2\) 的全部局部极值。
\end{enumerate}
\end{problemBox}

\problemAnchor{MATH1-CALC-0021}
\begin{problemBox}
\begin{problemMeta}
\textbf{题目编号：} \problemRef{MATH1-CALC-0021} \quad
\textbf{答案：} \answerRef{MATH1-CALC-0021} \quad
\textbf{来源：} Codex 原创 / K072 深度讲义配套练习（非真题）
\end{problemMeta}
\begin{enumerate}
  \item[\textbf{A}] 求 \(f(x)=x^2-2x\) 在 \([0,3]\) 上的最大值和最小值。
  \item[\textbf{B}] 求 \(f(x)=|x|-x\) 在 \([-2,3]\) 上的最大值和最小值，并写出全部取得点。
  \item[\textbf{C}] 求 \(f(x)=xe^{-x}\) 在 \([0,3]\) 上的最大值和最小值。
  \item[\textbf{M}] 讨论 \(a\in\mathbb R\) 时 \(f_a(x)=(x-a)^2\) 在 \([0,2]\) 上的最大值及其取得点。
\end{enumerate}
\end{problemBox}

\problemAnchor{MATH1-CALC-0022}
\begin{problemBox}
\begin{problemMeta}
\textbf{题目编号：} \problemRef{MATH1-CALC-0022} \quad
\textbf{答案：} \answerRef{MATH1-CALC-0022} \quad
\textbf{来源：} Codex 原创 / K073 深度讲义配套练习（非真题）
\end{problemMeta}
\begin{enumerate}
  \item[\textbf{A}] 求 \(f(x)=x^3+x^2\) 的凹凸区间和拐点。
  \item[\textbf{B}] 说明 \(f(x)=x^6+x^4\) 的 \(x=0\) 为什么不是拐点。
  \item[\textbf{C}] 求 \(f(x)=x-\ln x\,(x>0)\) 的凹凸区间和拐点。
  \item[\textbf{M}] 求 \(f(x)=x^4-2x^2\) 的凹凸区间与全部拐点。
\end{enumerate}
\end{problemBox}

\problemAnchor{MATH1-CALC-0023}
\begin{problemBox}
\begin{problemMeta}
\textbf{题目编号：} \problemRef{MATH1-CALC-0023} \quad
\textbf{答案：} \answerRef{MATH1-CALC-0023} \quad
\textbf{来源：} Codex 原创 / K074 深度讲义配套练习（非真题）
\end{problemMeta}
\begin{enumerate}
  \item[\textbf{A}] 求 \(f(x)=(x^2+1)/(x-2)\) 的全部渐近线。
  \item[\textbf{B}] 求 \(f(x)=\sqrt{4x^2+1}+2x\) 在 \(x\to\pm\infty\) 时的渐近线。
  \item[\textbf{C}] 判断 \(f(x)=1/(e^x-1)\) 的全部渐近线。
  \item[\textbf{M}] 求 \(f(x)=\sqrt{x^2+x+1}\) 在正、负无穷两端的斜渐近线。
\end{enumerate}
\end{problemBox}

\problemAnchor{MATH1-CALC-0024}
\begin{problemBox}
\begin{problemMeta}
\textbf{题目编号：} \problemRef{MATH1-CALC-0024} \quad
\textbf{答案：} \answerRef{MATH1-CALC-0024} \quad
\textbf{来源：} Codex 原创 / K075 深度讲义配套练习（非真题）
\end{problemMeta}
\begin{enumerate}
  \item[\textbf{A}] 根据零点、单调性、极值与拐点描绘 \(f(x)=x^4-4x^2\) 的草图。
  \item[\textbf{B}] 描绘 \(f(x)=x^2/(1+x^2)\) 的草图，并标出水平渐近线、极值点与拐点。
  \item[\textbf{C}] 描绘 \(f(x)=(\ln x)^2/x\,(x>0)\) 的草图，并标出两条渐近线、零点与极值点。
  \item[\textbf{M}] 描绘 \(f(x)=xe^{-x^2}\) 的草图。记 \(m=1/\sqrt{2e}\)，并据图说明方程 \(xe^{-x^2}=c\) 在 \(|c|>m,\ |c|=m,\ 0<|c|<m,\ c=0\) 时的实根个数。
\end{enumerate}
\end{problemBox}

\problemAnchor{MATH1-CALC-0025}
\begin{problemBox}
\begin{problemMeta}
\textbf{题目编号：} \problemRef{MATH1-CALC-0025} \quad
\textbf{答案：} \answerRef{MATH1-CALC-0025} \quad
\textbf{来源：} Codex 原创 / K076 深度讲义配套练习（非真题）
\end{problemMeta}
\begin{enumerate}
  \item[\textbf{A}] 求曲线 \(y=x^3\) 在 \(x=1\) 处的曲率。
  \item[\textbf{B}] 求曲线 \(y=\ln x\) 在 \(x=1\) 处的曲率。
  \item[\textbf{C}] 用参数公式求摆线 \(x=t-\sin t,\ y=1-\cos t\) 在 \(t=\pi\) 处的曲率。
  \item[\textbf{M}] 对参数曲线 \(x=t^2,y=t^3\)，求 \(t=1\) 处的曲率，并在 \(t>0\) 上消参后用显函数公式复核。
\end{enumerate}
\end{problemBox}

\problemAnchor{MATH1-CALC-0026}
\begin{problemBox}
\begin{problemMeta}
\textbf{题目编号：} \problemRef{MATH1-CALC-0026} \quad
\textbf{答案：} \answerRef{MATH1-CALC-0026} \quad
\textbf{来源：} Codex 原创 / K077 深度讲义配套练习（非真题）
\end{problemMeta}
\begin{enumerate}
  \item[\textbf{A}] 某正则曲线一点的曲率为 \(3/2\)，求该点的曲率半径。
  \item[\textbf{B}] 求曲线 \(y=e^x\) 在 \(x=0\) 处的曲率圆方程。
  \item[\textbf{C}] 求悬链线 \(y=a\cosh(x/a)\,(a>0)\) 在 \(x=0\) 处的曲率半径。
  \item[\textbf{M}] 说明曲线 \(y=\sin x\) 在原点为何没有普通有限曲率圆，并指出“曲率为零”时半径的正确退化含义。
\end{enumerate}
\end{problemBox}

\problemAnchor{MATH1-CALC-0027}
\begin{problemBox}
\begin{problemMeta}
\textbf{题目编号：} \problemRef{MATH1-CALC-0027} \quad
\textbf{答案：} \answerRef{MATH1-CALC-0027} \quad
\textbf{来源：} Codex 原创 / K078 深度讲义配套练习（非真题）
\end{problemMeta}
\begin{enumerate}
  \item[\textbf{A}] 判断并改正下列说法：“渐屈线属于本节点统考主线；Newton 一步给出的数就是精确根；数值迭代可以替代根的存在唯一性证明。”
  \item[\textbf{B}] 对抛物线 \(y=x^2\)，以横坐标 \(t\) 参数化其曲率中心，写出渐屈线的参数方程。
  \item[\textbf{C}] 对 \(f(x)=x^3-2\)，从 \(x_0=1\) 出发作一次 Newton 迭代，并说明所得数是精确值还是近似值。
  \item[\textbf{M}] 不使用数值迭代，证明方程 \(x^3+x-1=0\) 在 \((0,1)\) 上恰有一根，并把根定位在长度不超过 \(1/4\) 的区间内。
\end{enumerate}
\end{problemBox}
